\scanpage{85}
\mspage{69}
% archive foliation: 84
\deleted{
Appliquons \deleted{\unsure{8}}7.2.1 \deleted{\unsure{4}} au groupe formel
$\widehat A$, dans le cas favorable
envisagé dans \deleted{\unsure{10}} 7.1 où le schéma
\deleted{\unsure{dit}} formel abélien $\widehat B$ est algébrisable,
\inserted{et où on suppose de plus \deleted{\unsure{que}} $T_0$ isotrivial.}
On trouve alors une \deleted{extension canonique du} schéma
\deleted{abélien B dont le}\inserted{en groupes $A^{\unsure{\natural|\sharp}}$ sur $S$},
\inserted{appelé le \uline{groupe de Raynaud} associé au schéma en groupes $A$ sur $S$,
tel que $A$}
complété formel est $\widehat B$, par le tore $T$ sur $S$ qui remonte $T_0$,
\deleted{\unsure{;}} que nous
appellerons l'\uline{extension de Raynaud} associée à $A$, et que nous noterons
$A^{\unsure{\natural|\sharp}}$ :
\[
(\deleted{\unsure{8}}7.2.2)\qquad
0\longrightarrow T\longrightarrow A^{\unsure{\natural|\sharp}}
\longrightarrow B\longrightarrow0.
\]
Elle est donc définie, à isomorphisme unique près, par le fait qu'elle
donne naissance à un système cohérent d'isomorphismes d'extensions
\[
(7.2.3)\qquad A^{\unsure{\natural|\sharp}}_n\simeq A_n
\]
(le deuxième membre muni de sa structure d'extension décrite dans (7.1.1)
i.e. qu'on a un isomorphisme de \inserted{schémas formels en} groupes
\deleted{formels}
\[
(7.\unsure{2|3}.4)\qquad (A^{\unsure{\natural|\sharp}})^\wedge\simeq\widehat A,
\]
respectant d'ailleurs les structures d'extensions.
\deleted{Il résulte de plus de 7.2.1}

7.2.5. Il résulte d'autre part \deleted{que la} immédiatement de 7.2.1 que la
suite exacte (7.2.2) dépend fonctoriellement du schéma
\inserted{en groupes} $A$ sur $S$ satisfaisant aux conditions envisagées dans 7.1,
et que sa formation commute à
tout changement de base par un homomorphisme $S'\longrightarrow S$
associé à un homomorphisme d'anneaux topologiques noethériens adiques
$V\longrightarrow V'$.
}
